\UseRawInputEncoding
\documentclass[12pt,a4paper]{article}

% Packages for formatting and graphics
\usepackage{amsmath}
\usepackage{graphicx}
\usepackage{hyperref}
\usepackage{float}
\hypersetup{
    colorlinks=true,
    linkcolor=black,
    urlcolor=blue,
    citecolor=blue
}
\usepackage{geometry}
\geometry{a4paper, margin=1in}

\usepackage{listings}
\usepackage{color}
\usepackage{ragged2e}

\definecolor{codegreen}{rgb}{0,0.6,0}
\definecolor{codegray}{rgb}{0.5,0.5,0.5}
\definecolor{codepurple}{rgb}{0.58,0,0.82}
\definecolor{backcolour}{rgb}{0.95,0.95,0.92}

\lstdefinestyle{mystyle}{
    backgroundcolor=\color{backcolour},   
    commentstyle=\color{codegreen},
    keywordstyle=\color{magenta},
    numberstyle=\tiny\color{codegray},
    stringstyle=\color{codepurple},
    basicstyle=\ttfamily\footnotesize,
    breakatwhitespace=false,         
    breaklines=true,                 
    captionpos=b,                    
    keepspaces=true,                 
    numbers=left,                    
    numbersep=5pt,                  
    showspaces=false,                
    showstringspaces=false,
    showtabs=false,                  
    tabsize=2
}


% Title page
\title{EN3150 Assignment 03: Simple Convolutional Neural Network to Perform Classification}
\author{230415H Mudaliarachchi N.S. \\ 230185B Fernando S.D.}
\date{\today}

\begin{document}

\begin{titlepage}
    \begin{center}
        \LARGE
        \textbf{Department of Electronic and Telecommunication Engineering\\
        University of Moratuwa}

        \vspace{0.8cm}
        \textbf{EN3150 - Pattern Recognition}

        \vfill
        \includegraphics[width=0.3\textwidth]{university.jpeg}
        \vfill

        Assignment 03 \\
        Simple convolutional neural network to perform classification

        \vfill

        \Large
        \textbf{Group: Botzilla}

        \vspace{0.5cm}

        \begin{tabular}{cc}
            230415H & Mudaliarachchi N.S. \\
            230185B & Fernando S.D.       \\
        \end{tabular}

        \vfill

        \normalsize
        GitHub repository: \\
        \url{https://github.com/SenadiF/Botzilla---EN3150-Assignment-03.git}

        \vspace{0.8cm}
        \small 05/09/2026
    \end{center}
\end{titlepage}
\newpage

% Table of Contents
\tableofcontents
\newpage

\listoffigures
\listoftables
\newpage


\section{Resource-Constrained CNN for Edge Image Classification}

\subsection{Dataset and Class Imbalance Handling}

All models in this report were trained and evaluated on the RealWaste
dataset, which contains 4,752 colour images of waste items in nine
categories. The dataset is not balanced: as shown in
Table~\ref{tab:class_distribution}, Plastic is the largest class with
921 images (19.4\%), while Textile Trash is the smallest with 318 images
(6.7\%), giving an imbalance ratio of about $2.9:1$.

\begin{table}[H]
\centering
\caption{Class distribution of the RealWaste dataset and the test split}
\label{tab:class_distribution}
\begin{tabular}{lccc}
\hline
\textbf{Class} & \textbf{Images} & \textbf{Proportion} & \textbf{Test Samples} \\
\hline
Cardboard & 461 & 9.7\% & 69 \\
Food Organics & 411 & 8.6\% & 62 \\
Glass & 420 & 8.8\% & 63 \\
Metal & 790 & 16.6\% & 118 \\
Miscellaneous Trash & 495 & 10.4\% & 75 \\
Paper & 500 & 10.5\% & 75 \\
Plastic & 921 & 19.4\% & 138 \\
Textile Trash & 318 & 6.7\% & 48 \\
Vegetation & 436 & 9.2\% & 65 \\
\hline
\textbf{Total} & \textbf{4,752} & \textbf{100\%} & \textbf{713} \\
\hline
\end{tabular}
\end{table}

If this imbalance is ignored, a model can reach a reasonable overall
accuracy while performing poorly on the smaller classes, because the
larger classes contribute more to the training loss. The imbalance was
handled at two stages.

\paragraph{Data splitting.} The dataset was divided into training
(70\%, 3,326 images), validation (15\%, 713 images) and test (15\%,
713 images) subsets using stratified splitting. Stratification keeps the
proportion of each class the same in all three subsets as in the full
dataset, so that every class is represented in the validation and test
sets, and the evaluation results are not distorted by an uneven split.
The same split was used for every model in this report.

\paragraph{Training.} During training, class weights were applied in the
cross-entropy loss so that errors on the smaller classes are penalized
more heavily. Balanced class weights were used, computed from the
training set as
\begin{equation}
w_c = \frac{N}{K\,n_c},
\end{equation}
where $N$ is the number of training images, $K=9$ is the number of
classes and $n_c$ is the number of training images in class $c$. A class
with fewer samples therefore receives a larger weight; for example,
Textile Trash receives a weight of about 1.66 and Plastic about 0.57.
% CHECK: confirm this matches the class-weight formula used in the code.

Because the classes are not equally sized, weighted precision, recall and
F1-score are reported in addition to accuracy, and class-wise results are
given for every model.


\subsection{Model A -- Standard CNN}

Model A is a standard convolutional neural network consisting of four
standard $2$D convolutional layers, with each convolutional layer followed
by a ReLU activation and a $2\times2$ max pooling layer. The input image size
is $64\times64$ pixels with three RGB channels. The number of filters is
increased progressively from 32 to 256 to allow the network to learn
increasingly complex visual features.

The architecture of Model A is summarized below.

\begin{table}[H]
\centering
\caption{Architecture of Model A}
\label{tab:model_a_architecture}
\begin{tabular}{lcccc}
\hline
\textbf{Layer} & \textbf{Kernel} & \textbf{Filters/Units} &
\textbf{Activation} & \textbf{Output Size} \\
\hline
Input & -- & -- & -- & $3\times64\times64$ \\
Conv2D-1 & $3\times3$ & 32 & ReLU & $32\times64\times64$ \\
MaxPool & $2\times2$ & -- & -- & $32\times32\times32$ \\
Conv2D-2 & $3\times3$ & 64 & ReLU & $64\times32\times32$ \\
MaxPool & $2\times2$ & -- & -- & $64\times16\times16$ \\
Conv2D-3 & $3\times3$ & 128 & ReLU & $128\times16\times16$ \\
MaxPool & $2\times2$ & -- & -- & $128\times8\times8$ \\
Conv2D-4 & $3\times3$ & 256 & ReLU & $256\times8\times8$ \\
MaxPool & $2\times2$ & -- & -- & $256\times4\times4$ \\
Flatten & -- & -- & -- & 4096 \\
Fully Connected-1 & -- & 256 & ReLU & 256 \\
Fully Connected-2 & -- & 9 & -- & 9 \\
\hline
\end{tabular}
\end{table}

Each convolutional layer uses a $3\times3$ kernel with padding of 1. The
max pooling layers use a $2\times2$ kernel with a stride of 2, reducing the
spatial dimensions by half after each convolutional block. After the fourth
pooling layer, the resulting feature map has dimensions
$256\times4\times4$. This is flattened into 4096 features and passed through
a fully connected layer with 256 units, followed by the final fully connected
layer with 9 output units corresponding to the nine waste categories.

\subsubsection{Trainable Parameter Calculation}

The number of trainable parameters for a convolutional layer is calculated
using

\begin{equation}
N_{\mathrm{conv}} =
(K_h \times K_w \times C_{\mathrm{in}} \times C_{\mathrm{out}})
+ C_{\mathrm{out}},
\end{equation}

where the final term represents the bias parameters.

For the four convolutional layers:

\begin{align*}
\mathrm{Conv1} &: (3\times3\times3\times32)+32 = 896,\\
\mathrm{Conv2} &: (3\times3\times32\times64)+64 = 18,496,\\
\mathrm{Conv3} &: (3\times3\times64\times128)+128 = 73,856,\\
\mathrm{Conv4} &: (3\times3\times128\times256)+256 = 295,168.
\end{align*}

The number of parameters in a fully connected layer is calculated as

\begin{equation}
N_{\mathrm{FC}} =
N_{\mathrm{in}}\times N_{\mathrm{out}} + N_{\mathrm{out}}.
\end{equation}

Therefore,

\begin{align*}
\mathrm{FC1} &: (4096\times256)+256 = 1,048,832,\\
\mathrm{FC2} &: (256\times9)+9 = 2,313.
\end{align*}

Hence, the total number of trainable parameters is

\begin{align*}
N_{\mathrm{total}} &=
896 + 18,496 + 73,856 + 295,168\\
&\quad + 1,048,832 + 2,313\\
&= \boxed{1,439,561}.
\end{align*}

The ReLU, max-pooling, dropout, and flattening operations do not contain
trainable parameters. Therefore, Model A contains a total of
\textbf{1,439,561 trainable parameters}.

\subsubsection{Hardware Aware Activation Function}

ReLU was selected as the activation function for the convolutional and
hidden fully connected layers. ReLU is defined as

\begin{equation}
f(x)=\max(0,x).
\end{equation}

ReLU was selected because it is computationally inexpensive and only
requires a comparison with zero. Unlike sigmoid or softmax activations,
ReLU does not require computationally expensive exponential operations.
This makes it more suitable for resource-constrained hardware and reduces
the computational overhead during training and inference.

The final fully connected layer produces nine output logits without an
explicit softmax activation. Since Cross-Entropy Loss is used during
training, the required normalization is handled internally by the loss
function. Therefore, applying a separate softmax operation in the model
would introduce unnecessary exponential computations.

% ---------------------------------------------------------------------
\subsection{Optimizer Selection \& Tuning -- Model A}

For Model A, three optimization approaches were evaluated: Adam, standard
Stochastic Gradient Descent (SGD), and SGD with Momentum. Each optimizer was
tested using learning rates of 0.01, 0.001, and 0.0001. To ensure a fair
comparison, all experiments used the same CNN architecture, dataset split,
preprocessing pipeline, batch size, and training duration of 20 epochs. A new
randomly initialized model was used for each experiment.

\begin{table}[H]
\centering
\caption{Optimizer and Learning Rate Comparison for Model A}
\label{tab:optimizer_comparison_a}
\begin{tabular}{lcc}
\hline
\textbf{Optimizer} & \textbf{Learning Rate} & \textbf{Test Accuracy} \\
\hline
Adam & 0.01 & 36.04\% \\
\textbf{Adam} & \textbf{0.001} & \textbf{69.14\%} \\
Adam & 0.0001 & 57.92\% \\
\hline
SGD & 0.01 & 43.90\% \\
SGD & 0.001 & 19.35\% \\
SGD & 0.0001 & 12.76\% \\
\hline
SGD + Momentum & 0.01 & 60.17\% \\
SGD + Momentum & 0.001 & 42.50\% \\
SGD + Momentum & 0.0001 & 19.35\% \\
\hline
\end{tabular}
\end{table}

The results demonstrate that the learning rate had a substantial effect on
the convergence and final performance of Model A. For Adam, a learning rate
of 0.01 resulted in a test accuracy of 36.04\%, while reducing it to 0.001
increased the test accuracy to 69.14\%. Further reducing the learning rate
to 0.0001 resulted in a lower accuracy of 57.92\%. This indicates that, for
the selected CNN architecture and the fixed 20-epoch training period, a
learning rate of 0.001 provided a more suitable update magnitude among the
tested Adam configurations.

Standard SGD produced lower performance across the tested learning rates,
with its highest test accuracy being 43.90\% at a learning rate of 0.01.
At learning rates of 0.001 and 0.0001, the model showed limited convergence
within the 20 training epochs, resulting in test accuracies of 19.35\% and
12.76\%, respectively.

SGD with Momentum was also evaluated using a momentum value of 0.9.
Momentum incorporates information from previous parameter updates into the
current update, helping the optimizer maintain movement in consistent
directions and potentially improving convergence. With a learning rate of
0.01, adding momentum increased the test accuracy from 43.90\% with
standard SGD to 60.17\%. At a learning rate of 0.001, SGD with Momentum
achieved 42.50\%, compared with 19.35\% for standard SGD.

Based on these experiments, Adam with a learning rate of 0.001 was selected
for Model A because it achieved the highest test accuracy among the
evaluated optimizer and learning-rate combinations.

% ---------------------------------------------------------------------
\subsection{Custom Model Training \& Evaluation -- Model A}

Model A was trained for 20 epochs using the Adam optimizer with a learning
rate of 0.001, which was selected during the optimizer tuning experiments.
The images were resized to $64\times64$ pixels, the stratified
training, validation, and test split was used, and a batch size of 64 was
used during training.

Because the weight initialization and the order of the mini-batches are
random, a single training run can give an optimistic or pessimistic
picture of the model. Model A was therefore trained five independent
times, each from a new random initialization with a different random
seed, while the data split, preprocessing and hyperparameters were kept
fixed. Each trained model was evaluated on the same test set of 713
images. Results are reported as the mean $\pm$ sample standard deviation
over the five runs,

\begin{equation}
\bar{x}=\frac{1}{n}\sum_{i=1}^{n}x_i, \qquad
s=\sqrt{\frac{1}{n-1}\sum_{i=1}^{n}\left(x_i-\bar{x}\right)^2},
\qquad n=5.
\end{equation}

Training and validation losses were recorded at the end of every epoch.
Figure~\ref{fig:model_a_loss} shows the mean loss across the five runs,
with the shaded region indicating $\pm1$ standard deviation.

\begin{figure}[H]
    \centering
    \includegraphics[width=0.85\linewidth]{ModelA_loss_mean_std.png}
    \caption{Training and validation loss of Model A, averaged over five
    independent runs. Shaded bands show $\pm1$ standard deviation.}
    \label{fig:model_a_loss}
\end{figure}

The training loss decreased steadily from about 2.05 to 0.67, and the
validation loss decreased from about 1.86 to 0.75, with a narrow spread
between runs. The validation loss fluctuates more than the training loss
in the final epochs, and the gap between the two begins to widen slightly,
suggesting that the model is close to the point where further training
without regularization would lead to overfitting.

The test results of the individual runs are listed in
Table~\ref{tab:model_a_runs}.

\begin{table}[H]
\centering
\caption{Test performance of Model A over five independent runs}
\label{tab:model_a_runs}
\resizebox{\textwidth}{!}{
\begin{tabular}{cccccc}
\hline
\textbf{Run} & \textbf{Accuracy (\%)} & \textbf{W. Precision} &
\textbf{W. Recall} & \textbf{W. F1} & \textbf{Train Time (s)} \\
\hline
1 & 72.93 & 0.7206 & 0.7293 & 0.7208 & 745.82 \\
2 & 71.39 & 0.7224 & 0.7139 & 0.7154 & 876.61 \\
3 & 69.71 & 0.7082 & 0.6971 & 0.6940 & 884.21 \\
4 & 69.99 & 0.6981 & 0.6999 & 0.6970 & 880.30 \\
5 & 72.79 & 0.7427 & 0.7279 & 0.7290 & 900.18 \\
\hline
\textbf{Mean $\pm$ SD} & \textbf{71.36 $\pm$ 1.51} & \textbf{0.7184 $\pm$ 0.0168}
& \textbf{0.7136 $\pm$ 0.0151} & \textbf{0.7112 $\pm$ 0.0152}
& \textbf{857.4 $\pm$ 63.0} \\
\hline
\end{tabular}
}
\end{table}

Over the five runs, Model A achieved a mean test accuracy of
\textbf{71.36 $\pm$ 1.51\%}, a mean weighted precision of
\textbf{71.84 $\pm$ 1.68\%}, a mean weighted recall of
\textbf{71.36 $\pm$ 1.51\%}, and a mean weighted F1-score of
\textbf{71.12 $\pm$ 1.52\%}. Test accuracy ranged from 69.71\% to
72.93\%, so the model's performance was consistent across random
initializations. The mean training time was $857.4\pm63.0$\,s for 20
epochs on a CPU, or about $42.9\pm3.2$\,s per epoch.

\begin{table}[H]
\centering
\caption{Overall Test Performance of Model A (mean $\pm$ SD, $n=5$)}
\label{tab:model_a_overall}
\begin{tabular}{lc}
\hline
\textbf{Metric} & \textbf{Model A} \\
\hline
Test Accuracy & 71.36 $\pm$ 1.51\% \\
Weighted Precision & 71.84 $\pm$ 1.68\% \\
Weighted Recall & 71.36 $\pm$ 1.51\% \\
Weighted F1-score & 71.12 $\pm$ 1.52\% \\
Training Time per Epoch & 42.9 $\pm$ 3.2\,s \\
Test Samples & 713 \\
\hline
\end{tabular}
\end{table}

Each of the five trained models was also evaluated class by class. The
class-wise precision, recall and F1-score, averaged over the five runs,
are given in Table~\ref{tab:model_a_classwise}.

\begin{table}[H]
\centering
\caption{Class-wise performance of Model A (mean $\pm$ SD, $n=5$)}
\label{tab:model_a_classwise}
\small
\begin{tabular}{lccc}
\hline
\textbf{Class} & \textbf{Precision} & \textbf{Recall} & \textbf{F1-score} \\
\hline
Cardboard & 0.757 $\pm$ 0.021 & 0.785 $\pm$ 0.038 & 0.771 $\pm$ 0.026 \\
Food Organics & 0.822 $\pm$ 0.045 & 0.755 $\pm$ 0.064 & 0.785 $\pm$ 0.029 \\
Glass & 0.773 $\pm$ 0.080 & 0.679 $\pm$ 0.052 & 0.719 $\pm$ 0.026 \\
Metal & 0.692 $\pm$ 0.029 & 0.731 $\pm$ 0.040 & 0.710 $\pm$ 0.014 \\
Miscellaneous Trash & 0.475 $\pm$ 0.036 & 0.453 $\pm$ 0.123 & 0.458 $\pm$ 0.071 \\
Paper & 0.795 $\pm$ 0.066 & 0.725 $\pm$ 0.053 & 0.756 $\pm$ 0.027 \\
Plastic & 0.681 $\pm$ 0.053 & 0.754 $\pm$ 0.067 & 0.712 $\pm$ 0.025 \\
Textile Trash & 0.683 $\pm$ 0.104 & 0.529 $\pm$ 0.054 & 0.590 $\pm$ 0.033 \\
Vegetation & 0.873 $\pm$ 0.050 & 0.938 $\pm$ 0.039 & 0.903 $\pm$ 0.021 \\
\hline
\textbf{Macro Average} & 0.728 $\pm$ 0.023 & 0.706 $\pm$ 0.018 & 0.711 $\pm$ 0.016 \\
\hline
\end{tabular}
\end{table}

Vegetation was the easiest class in every run (mean F1 0.903), followed
by Food Organics and Cardboard. Miscellaneous Trash was consistently the
hardest class (mean F1 0.458), followed by Textile Trash (0.590). These
two classes also had the largest run-to-run variation: Miscellaneous
Trash recall ranged from 0.32 to 0.65 across the five runs, which shows
that the model's handling of this heterogeneous category depends heavily
on the random initialization.

To show where the errors occur, the confusion matrices of the five runs
were averaged. Since every run was evaluated on the same 713 test
images, the matrices can be averaged cell by cell. The result is shown
in Figure~\ref{fig:model_a_confusion}. Rows represent the actual classes
and columns the predicted classes. The left panel shows the mean count
per cell, and the right panel normalizes each row by the number of
samples in that class, so the diagonal gives the mean recall of each
class and classes of different sizes can be compared fairly; the value
below each percentage is its standard deviation across runs. The
confusion matrices of the individual runs are given in
Section~\ref{sec:runs_a}.

\begin{figure}[H]
    \centering
    \includegraphics[width=\linewidth]{ModelA_confusion_avg.png}
    \caption{Confusion matrix of Model A on the test set, averaged over
    five runs: mean counts (left) and mean percentage of each true class
    with standard deviation (right).}
    \label{fig:model_a_confusion}
\end{figure}

% TODO: add 1-2 sentences on the largest average confusions, using the
% "Largest average confusions" lines printed by average_confusion.py
% (e.g. "On average, x% of Plastic samples were predicted as Metal").

Individual runs can differ noticeably from this average. For example,
Miscellaneous Trash recall was 65.3\% (49 of 75) in Run 2
(Figure~\ref{fig:model_a_run2}) but only 32.0\% (24 of 75) in Run 1,
against a five-run mean of 45.3\%. A single confusion matrix can
therefore over- or under-state the difficulty of individual classes,
which is why the averaged results are used for the overall conclusions.


\subsubsection{Confusion Matrices of Individual Runs -- Model A}
\label{sec:runs_a}

The confusion matrices of the five individual runs of Model A are shown
below. Each figure shows raw counts (left) and the percentage of each
true class (right).

\begin{figure}[H]
    \centering
    \includegraphics[width=\linewidth]{ModelA_confusion_run1.png}
    \caption{Model A, Run 1 (test accuracy 72.93\%).}
\end{figure}

\begin{figure}[H]
    \centering
    \includegraphics[width=\linewidth]{ModelA_confusion_run2.png}
    \caption{Model A, Run 2 (test accuracy 71.39\%).}
    \label{fig:model_a_run2}
\end{figure}

\begin{figure}[H]
    \centering
    \includegraphics[width=\linewidth]{ModelA_confusion_run3.png}
    \caption{Model A, Run 3 (test accuracy 69.71\%).}
\end{figure}

\begin{figure}[H]
    \centering
    \includegraphics[width=\linewidth]{ModelA_confusion_run4.png}
    \caption{Model A, Run 4 (test accuracy 69.99\%).}
\end{figure}

\begin{figure}[H]
    \centering
    \includegraphics[width=\linewidth]{ModelA_confusion_run5.png}
    \caption{Model A, Run 5 (test accuracy 72.79\%).}
\end{figure}

\subsection{Model B -- Lightweight CNN}

Model B was designed as a lightweight convolutional neural network for
waste image classification under resource constrained hardware conditions.
Unlike Model A, which uses standard convolutional layers, Model B uses
depthwise separable convolutions to reduce the number of trainable
parameters and computational requirements.

The model accepts RGB images resized to $64\times64$ pixels and classifies
them into the nine waste categories in the RealWaste dataset. It consists
of three depthwise-separable convolutional blocks, each followed by a
$2\times2$ max pooling layer. The number of output channels increases
from 32 to 64 and then to 128. Adaptive average pooling is subsequently
used to reduce the spatial dimensions to $1\times1$. The resulting
128 features are passed through a fully connected layer with 256 units
and a final output layer with nine units.

Batch Normalization is applied after the depthwise and pointwise
convolutions, followed by ReLU activation. Dropout with a probability
of 0.3 is used in the classifier to help reduce overfitting.

The architecture of Model B is summarized in
Table~\ref{tab:model_b_architecture}.

\begin{table}[H]
\centering
\caption{Architecture of Model B}
\label{tab:model_b_architecture}
\small
\begin{tabular}{lcc}
\hline
\textbf{Layer} & \textbf{Output Channels/Units} &
\textbf{Output Size} \\
\hline
Input & 3 & $3\times64\times64$ \\
Depthwise Conv2D-1 & 3 & $3\times64\times64$ \\
Pointwise Conv2D-1 & 32 & $32\times64\times64$ \\
MaxPool & -- & $32\times32\times32$ \\
Depthwise Conv2D-2 & 32 & $32\times32\times32$ \\
Pointwise Conv2D-2 & 64 & $64\times32\times32$ \\
MaxPool & -- & $64\times16\times16$ \\
Depthwise Conv2D-3 & 64 & $64\times16\times16$ \\
Pointwise Conv2D-3 & 128 & $128\times16\times16$ \\
MaxPool & -- & $128\times8\times8$ \\
Adaptive Average Pooling & 128 & $128\times1\times1$ \\
Flatten & -- & 128 \\
Dropout & -- & 128 \\
Fully Connected-1 & 256 & 256 \\
ReLU and Dropout & -- & 256 \\
Fully Connected-2 & 9 & 9 \\
\hline
\end{tabular}
\end{table}

\subsubsection{Trainable Parameter Calculation}

A depthwise-separable convolution separates a standard convolution
into two operations: depthwise convolution and pointwise convolution.

The depthwise convolution applies a separate $3\times3$ filter to each
input channel. Its number of convolutional parameters is

\begin{equation}
N_{\mathrm{depthwise}} = K_h K_w C_{\mathrm{in}},
\end{equation}

where $K_h$ and $K_w$ are the kernel dimensions and
$C_{\mathrm{in}}$ is the number of input channels. Bias terms are
not used in the convolutional layers.

The pointwise convolution uses a $1\times1$ kernel to combine information
across channels. Its number of parameters is

\begin{equation}
N_{\mathrm{pointwise}} =
C_{\mathrm{in}}C_{\mathrm{out}}.
\end{equation}

Batch Normalization introduces two trainable parameters per channel:
a scale parameter and a shift parameter. Therefore, the trainable
parameter count for each depthwise-separable block is

\begin{equation}
N_{\mathrm{block}} =
K_hK_wC_{\mathrm{in}} + 2C_{\mathrm{in}}
+ C_{\mathrm{in}}C_{\mathrm{out}} + 2C_{\mathrm{out}}.
\end{equation}

For the three convolutional blocks, the parameter counts are

\begin{align*}
\mathrm{Block\ 1}
&= (3\times3\times3)+(2\times3)\\
&\quad +(3\times32)+(2\times32)\\
&=193,\\[4pt]
\mathrm{Block\ 2}
&= (3\times3\times32)+(2\times32)\\
&\quad +(32\times64)+(2\times64)\\
&=2,528,\\[4pt]
\mathrm{Block\ 3}
&= (3\times3\times64)+(2\times64)\\
&\quad +(64\times128)+(2\times128)\\
&=9,152.
\end{align*}

The fully connected layers contain

\begin{align*}
\mathrm{FC1} &= (128\times256)+256 = 33,024,\\
\mathrm{FC2} &= (256\times9)+9 = 2,313.
\end{align*}

Adaptive average pooling, max pooling, ReLU, and dropout do not
introduce trainable parameters. Therefore, the total number of
trainable parameters is

\begin{align*}
N_{\mathrm{total}}
&=193+2,528+9,152\\
&\quad +33,024+2,313\\
&=\boxed{47,210}.
\end{align*}

Thus, Model B contains 47,210 trainable parameters.

\subsubsection{Hardware Aware Activation Function}

ReLU was selected as the activation function in the convolutional
blocks and the hidden fully connected layer. It is defined as

\begin{equation}
f(x)=\max(0,x).
\end{equation}

ReLU is computationally inexpensive because it only requires comparing
each input with zero. This avoids the exponential operations required
by sigmoid and softmax activations and makes ReLU suitable for
resource constrained devices.

Batch Normalization is used to normalize intermediate feature
activations during training, while dropout randomly deactivates a
fraction of hidden units during training to help reduce overfitting.

The final fully connected layer produces nine output logits without
an explicit softmax activation. Cross Entropy Loss is used during
training, and it internally performs the required log softmax
computation. Therefore, a separate softmax layer is not required
during training.

% ---------------------------------------------------------------------

\subsection{Custom Model Training \& Evaluation -- Model B}

Model B was trained for 20 epochs using the Adam optimizer with a
learning rate of $0.001$, using the same $64\times64$ input size,
stratified data split, and batch size of 64 as Model A. Following the
same protocol as Model A, Model B was trained five independent times
with random seeds 1001--1005, and each run was evaluated on the same
713-image test set.

The mean training and validation loss over the five runs is shown in
Figure~\ref{fig:model_b_loss}, with the shaded region indicating
$\pm1$ standard deviation.

\begin{figure}[H]
    \centering
    \includegraphics[width=0.85\linewidth]{ModelB_loss_mean_std.png}
    \caption{Training and validation loss of Model B, averaged over five
    independent runs. Shaded bands show $\pm1$ standard deviation.}
    \label{fig:model_b_loss}
\end{figure}

The mean training loss decreased from about 2.02 to 1.04, and the mean
validation loss decreased from about 1.85 to 1.04. Two observations
stand out when compared with Model A. First, the final losses are
considerably higher than Model A's (about 0.67 training and 0.75
validation), and the training and validation curves stay close together
throughout. This indicates that Model B is underfitting rather than
overfitting: its limited capacity, together with dropout, prevents it
from fitting the training data closely, and the training loss is still
decreasing at epoch 20, so additional epochs would likely improve it.
Second, the validation loss is much noisier and its standard deviation
band is wider than Model A's. Three of the five runs showed a sharp,
temporary increase in validation loss at epoch 15, which produces the
spike in the mean curve. The final validation loss also varied more
between runs ($1.04\pm0.07$), which is consistent with the larger spread
in Model B's test accuracy.

The test results of the individual runs are listed in
Table~\ref{tab:model_b_runs}.

\begin{table}[H]
\centering
\caption{Test performance of Model B over five independent runs}
\label{tab:model_b_runs}
\resizebox{\textwidth}{!}{
\begin{tabular}{ccccccc}
\hline
\textbf{Run} & \textbf{Seed} & \textbf{Accuracy (\%)} & \textbf{W. Precision} &
\textbf{W. Recall} & \textbf{W. F1} & \textbf{Train Time (s)} \\
\hline
1 & 1001 & 67.74 & 0.6825 & 0.6774 & 0.6730 & 830.97 \\
2 & 1002 & 63.53 & 0.6515 & 0.6353 & 0.6225 & 813.51 \\
3 & 1003 & 57.78 & 0.5961 & 0.5778 & 0.5698 & 745.62 \\
4 & 1004 & 64.38 & 0.6446 & 0.6438 & 0.6216 & 831.36 \\
5 & 1005 & 62.13 & 0.6237 & 0.6213 & 0.6007 & 842.35 \\
\hline
\multicolumn{2}{c}{\textbf{Mean $\pm$ SD}} & \textbf{63.11 $\pm$ 3.62}
& \textbf{0.6397 $\pm$ 0.0322} & \textbf{0.6311 $\pm$ 0.0362}
& \textbf{0.6175 $\pm$ 0.0377} & \textbf{812.8 $\pm$ 38.9} \\
\hline
\end{tabular}
}
\end{table}

Over the five runs, Model B achieved a mean test accuracy of
\textbf{63.11 $\pm$ 3.62\%}, a mean weighted precision of
\textbf{63.97 $\pm$ 3.22\%}, a mean weighted recall of
\textbf{63.11 $\pm$ 3.62\%}, and a mean weighted F1-score of
\textbf{61.75 $\pm$ 3.77\%}. Test accuracy ranged from 57.78\% to
67.74\%. The standard deviation is more than twice that of Model A,
which indicates that the small model is more sensitive to its random
initialization and that 20 epochs is not always enough for it to converge
to a good solution.

\begin{table}[H]
\centering
\caption{Overall Test Performance of Model B (mean $\pm$ SD, $n=5$)}
\label{tab:model_b_overall}
\begin{tabular}{lc}
\hline
\textbf{Metric} & \textbf{Model B} \\
\hline
Test Accuracy & 63.11 $\pm$ 3.62\% \\
Weighted Precision & 63.97 $\pm$ 3.22\% \\
Weighted Recall & 63.11 $\pm$ 3.62\% \\
Weighted F1-score & 61.75 $\pm$ 3.77\% \\
Training Time per Epoch & 40.6 $\pm$ 1.9\,s \\
Test Samples & 713 \\
\hline
\end{tabular}
\end{table}

The class-wise precision, recall and F1-score of Model B, averaged over
the five runs, are given in Table~\ref{tab:model_b_classwise}.

\begin{table}[H]
\centering
\caption{Class-wise performance of Model B (mean $\pm$ SD, $n=5$)}
\label{tab:model_b_classwise}
\small
\begin{tabular}{lccc}
\hline
\textbf{Class} & \textbf{Precision} & \textbf{Recall} & \textbf{F1-score} \\
\hline
Cardboard & 0.809 $\pm$ 0.053 & 0.786 $\pm$ 0.052 & 0.795 $\pm$ 0.020 \\
Food Organics & 0.768 $\pm$ 0.043 & 0.790 $\pm$ 0.127 & 0.777 $\pm$ 0.081 \\
Glass & 0.790 $\pm$ 0.044 & 0.492 $\pm$ 0.050 & 0.604 $\pm$ 0.037 \\
Metal & 0.581 $\pm$ 0.076 & 0.649 $\pm$ 0.165 & 0.594 $\pm$ 0.064 \\
Miscellaneous Trash & 0.404 $\pm$ 0.059 & 0.277 $\pm$ 0.125 & 0.313 $\pm$ 0.096 \\
Paper & 0.613 $\pm$ 0.101 & 0.613 $\pm$ 0.030 & 0.609 $\pm$ 0.054 \\
Plastic & 0.574 $\pm$ 0.066 & 0.730 $\pm$ 0.040 & 0.639 $\pm$ 0.025 \\
Textile Trash & 0.509 $\pm$ 0.149 & 0.246 $\pm$ 0.101 & 0.311 $\pm$ 0.085 \\
Vegetation & 0.837 $\pm$ 0.078 & 0.920 $\pm$ 0.048 & 0.873 $\pm$ 0.030 \\
\hline
\textbf{Macro Average} & 0.654 $\pm$ 0.040 & 0.612 $\pm$ 0.035 & 0.613 $\pm$ 0.038 \\
\hline
\end{tabular}
\end{table}

As for Model A, the confusion matrices of the five runs were averaged
cell by cell; the result is shown in Figure~\ref{fig:model_b_confusion}.
The individual runs are given in Section~\ref{sec:runs_b}.

\begin{figure}[H]
    \centering
    \includegraphics[width=\linewidth]{ModelB_confusion_avg.png}
    \caption{Confusion matrix of Model B on the test set, averaged over
    five runs: mean counts (left) and mean percentage of each true class
    with standard deviation (right).}
    \label{fig:model_b_confusion}
\end{figure}

% TODO: add 1-2 sentences on the largest average confusions, using the
% "Largest average confusions" lines printed by average_confusion.py.

Across the five runs, Model B performed best on Vegetation, Cardboard
and Food Organics (mean F1 0.873, 0.795 and 0.777). Its weakest classes
were the same as Model A's, but with much lower scores: Miscellaneous
Trash (mean F1 0.313) and Textile Trash (0.311), with mean recalls of
only 0.277 and 0.246. Glass also had low recall (0.492) despite high
precision (0.790), meaning that the model rarely predicted Glass
wrongly but missed about half of the actual Glass samples. Model B's
class-level results also varied much more between runs than Model A's:
Metal recall ranged from 0.37 to 0.77 and Miscellaneous Trash recall
from 0.13 to 0.44.

The individual runs show the same pattern. In Run 2
(Figure~\ref{fig:model_b_run2}), for example, Cardboard
(57 of 69, 82.6\%) and Food Organics (55 of 62, 88.7\%) were classified
well, whereas only 7 of 48 Textile Trash samples (14.6\%), 25 of 75
Miscellaneous Trash samples (33.3\%) and 30 of 63 Glass samples
(47.6\%) were correct.

These results indicate that the lightweight architecture learned
useful visual features but had difficulty distinguishing some
visually similar waste categories. The relatively small number of
trainable parameters reduces model complexity, but it may also
limit the capacity to learn more discriminative features.

\subsubsection{Confusion Matrices of Individual Runs -- Model B}
\label{sec:runs_b}

The confusion matrices of the five individual runs of Model B are shown
below. Each figure shows raw counts (left) and the percentage of each
true class (right).

\begin{figure}[H]
    \centering
    \includegraphics[width=\linewidth]{ModelB_confusion_run1.png}
    \caption{Model B, Run 1, seed 1001 (test accuracy 67.74\%).}
\end{figure}

\begin{figure}[H]
    \centering
    \includegraphics[width=\linewidth]{ModelB_confusion_run2.png}
    \caption{Model B, Run 2, seed 1002 (test accuracy 63.53\%).}
    \label{fig:model_b_run2}
\end{figure}

\begin{figure}[H]
    \centering
    \includegraphics[width=\linewidth]{ModelB_confusion_run3.png}
    \caption{Model B, Run 3, seed 1003 (test accuracy 57.78\%).}
\end{figure}

\begin{figure}[H]
    \centering
    \includegraphics[width=\linewidth]{ModelB_confusion_run4.png}
    \caption{Model B, Run 4, seed 1004 (test accuracy 64.38\%).}
\end{figure}

\begin{figure}[H]
    \centering
    \includegraphics[width=\linewidth]{ModelB_confusion_run5.png}
    \caption{Model B, Run 5, seed 1005 (test accuracy 62.13\%).}
\end{figure}

% =====================================================================
% COMPARISON A vs B
% =====================================================================
\subsection{Comparison of Model A and Model B}

Model A and Model B were compared in terms of total parameter count,
estimated model size on disk, training time per epoch, and test
accuracy. The estimated model size assumes that each parameter is stored
as a 32-bit (4-byte) floating-point value, i.e.\
$\text{size} \approx 4 \times N_{\mathrm{params}}$ bytes. The number of
multiply--accumulate operations (MACs) per image was also calculated
analytically from the layer dimensions as an indicator of computational
cost.

Accuracy, F1-score and training time are reported as mean $\pm$
standard deviation over five independent runs of each model.

\begin{table}[H]
\centering
\caption{Comparison of Model A and Model B (performance metrics: mean $\pm$ SD, $n=5$)}
\label{tab:model_ab_comparison}
\begin{tabular}{lcc}
\hline
\textbf{Metric} & \textbf{Model A} & \textbf{Model B} \\
\hline
Convolution Type & Standard & Depthwise-separable \\
Total Trainable Parameters & 1,439,561 & 47,210 \\
Estimated Model Size on Disk (KB) & 5,623.3 & 184.4 \\
Approx. MACs per Image & 61.2 M & 5.18 M \\
Training Time per Epoch (s, CPU) & 42.9 $\pm$ 3.2 & 40.6 $\pm$ 1.9 \\
Test Accuracy & 71.36 $\pm$ 1.51\% & 63.11 $\pm$ 3.62\% \\
Weighted F1-score & 71.12 $\pm$ 1.52\% & 61.75 $\pm$ 3.77\% \\
\hline
\end{tabular}
\end{table}

Although Model B needs about $11.8\times$ fewer MACs per image, its
training time per epoch was only slightly lower than that of Model A.
On a CPU, depthwise convolutions are memory-bound rather than
compute-bound and are less well optimized in standard deep learning
libraries than dense convolutions, and training time also includes data
loading and image resizing, which is identical for both models. The
reduction in MACs is therefore not fully reflected in wall-clock training
time on this hardware.

\paragraph{Trade-offs of moving to depthwise-separable convolutions.}
A standard $3\times3$ convolution learns spatial and cross-channel
patterns jointly using $9C_{\mathrm{in}}C_{\mathrm{out}}$ weights. A
depthwise-separable convolution factorizes this into a per-channel
$3\times3$ spatial filter ($9C_{\mathrm{in}}$ weights) followed by a
$1\times1$ channel-mixing filter ($C_{\mathrm{in}}C_{\mathrm{out}}$
weights). For layers with many output channels, this reduces the
parameter and computation cost by a factor close to
$1/C_{\mathrm{out}} + 1/9$, i.e.\ roughly $8$--$9\times$ per layer.

Model B combines this factorization with two further design choices:
it uses three convolutional blocks instead of four, and it replaces the
large flatten-plus-dense layer of Model A ($4096\times256$, which alone
accounts for 1,048,832 of Model A's 1,439,561 parameters) with adaptive
average pooling. As a result, Model B has approximately 96.72\% fewer
trainable parameters than Model A ($30.5\times$ smaller), its estimated
storage drops from about 5.5\,MB to about 184\,KB, and its computational
cost falls from approximately 61.2 million to 5.18 million MACs per image
(about $11.8\times$ fewer).

However, this compression came with a loss in accuracy. Model A achieved
a mean test accuracy of $71.36\pm1.51\%$ whereas Model B achieved
$63.11\pm3.62\%$, a drop of 8.25 percentage points, and the mean weighted
F1-score fell from 71.12\% to 61.75\%. The gap is several times larger
than the run-to-run standard deviation of either model, and the
accuracy ranges of the two models (69.71--72.93\% and 57.78--67.74\%) do
not overlap, so the difference is not an artifact of a particular
initialization. Model B was also less stable, with a standard deviation
more than twice that of Model A. The drop was concentrated in the harder,
visually diverse classes: averaged over five runs, Miscellaneous Trash
recall fell from 0.45 to 0.28 and Textile Trash recall from 0.53 to
0.25, while Glass recall fell from 0.68 to 0.49. In contrast, the easy
classes changed little (Cardboard recall was 0.79 for both models and
Vegetation 0.94 versus 0.92). Textile Trash is also the smallest class
in the dataset (6.7\%), so even with class-weighted training it provides
the fewest examples to learn from, and the smaller Model B was less able
to make use of them. Because each depthwise filter only
sees one channel, and the network has fewer layers and channels, Model B
has less representational capacity to learn the fine-grained texture and
colour combinations that separate these classes.

In practice, Model A is the better choice if accuracy is the priority and
the target hardware has several megabytes of memory available. Model B is
far better suited to tightly constrained edge devices, where its sub-200\,KB footprint and low compute
requirement make deployment feasible.


\newpage
\section{Compare with Lightweight State-of-the-Art Networks}

% ---------------------------------------------------------------------
\subsection{SOTA Fine-Tuning \& Evaluation}

Two pre-trained lightweight convolutional neural network architectures,
EfficientNet-B0 and MobileNetV2, were selected for evaluation. Both
architectures are designed to provide efficient image classification and
are suitable for resource-constrained and mobile/edge computing
environments. The same training, validation, and test splits were used
for both models to ensure a consistent comparison. The dataset was
divided into 70\% training, 15\% validation, and 15\% testing subsets
using the same stratified split. Both models used ImageNet pretrained
weights, and input images were resized to $224\times224$ pixels.

Both models followed the same two-stage transfer-learning strategy: an
initial stage in which the pretrained backbone was frozen and only a new
classification layer was trained, followed by low-learning-rate
fine-tuning of the deeper feature layers.

\subsubsection{EfficientNet-B0}

EfficientNet-B0 was initialized using ImageNet-pretrained weights. Since
the original model was trained to classify 1000 ImageNet classes, its
final classification layer was replaced with a new fully connected layer
producing 9 outputs, corresponding to the nine waste categories in the
RealWaste dataset.

Initially, the EfficientNet-B0 feature extractor was frozen and only the
newly added classifier was trained. The model was trained for 10 epochs
using the Adam optimizer with a learning rate of $0.001$. At the end of
this stage, the model achieved a validation accuracy of $83.03\%$.

The deeper feature layers were then fine-tuned by unfreezing feature
blocks 6--8 while keeping the earlier feature blocks frozen. The learning
rate was reduced to $0.0001$, and the model was fine-tuned for an
additional 5 epochs. Validation accuracy increased to $89.90\%$ at the end
of this stage.

To investigate whether additional fine-tuning could improve the model, a
further 5 epochs were performed using the same learning rate and
trainable layers. The resulting model achieved a test loss of $0.3521$
and a test accuracy of $88.92\%$.

\begin{table}[H]
    \centering
    \caption{EfficientNet-B0 classification report on the RealWaste test set.}
    \label{tab:efficientnet_classification}
    \begin{tabular}{lcccc}
        \hline
        Class & Precision & Recall & F1-score & Support \\
        \hline
        Cardboard & 0.92 & 0.94 & 0.93 & 69 \\
        Food Organics & 0.93 & 0.90 & 0.92 & 62 \\
        Glass & 0.93 & 0.84 & 0.88 & 63 \\
        Metal & 0.87 & 0.89 & 0.88 & 118 \\
        Miscellaneous Trash & 0.79 & 0.80 & 0.79 & 75 \\
        Paper & 0.94 & 0.88 & 0.91 & 75 \\
        Plastic & 0.85 & 0.89 & 0.87 & 138 \\
        Textile Trash & 0.90 & 0.92 & 0.91 & 48 \\
        Vegetation & 0.97 & 0.95 & 0.96 & 65 \\
        \hline
        Accuracy & & & 0.89 & 713 \\
        Macro Avg. & 0.90 & 0.89 & 0.89 & 713 \\
        Weighted Avg. & 0.89 & 0.89 & 0.89 & 713 \\
        \hline
    \end{tabular}
\end{table}

The confusion matrix for EfficientNet-B0 is shown in Table~\ref{tab:efficientnet_confusion}. 

\begin{table}[H]
    \centering
    \caption{EfficientNet-B0 confusion matrix on the RealWaste test set.}
    \label{tab:efficientnet_confusion}
    \resizebox{\textwidth}{!}{
    \begin{tabular}{lccccccccc}
        \hline
        & Cardboard & Food Org. & Glass & Metal & Misc. Trash & Paper & Plastic & Textile & Vegetation \\
        \hline
        Cardboard
        & 65 & 0 & 0 & 0 & 0 & 1 & 3 & 0 & 0 \\
        Food Organics
        & 0 & 56 & 0 & 0 & 4 & 0 & 0 & 1 & 1 \\
        Glass
        & 0 & 0 & 53 & 3 & 2 & 1 & 4 & 0 & 0 \\
        Metal
        & 1 & 0 & 0 & 105 & 4 & 1 & 7 & 0 & 0 \\
        Misc. Trash
        & 1 & 1 & 1 & 3 & 60 & 1 & 5 & 3 & 0 \\
        Paper
        & 4 & 1 & 0 & 1 & 0 & 66 & 2 & 1 & 0 \\
        Plastic
        & 0 & 0 & 3 & 8 & 3 & 0 & 123 & 0 & 1 \\
        Textile Trash
        & 0 & 0 & 0 & 1 & 2 & 0 & 1 & 44 & 0 \\
        Vegetation
        & 0 & 2 & 0 & 0 & 1 & 0 & 0 & 0 & 62 \\
        \hline
    \end{tabular}
    }
\end{table}

The EfficientNet-B0 model correctly classified most samples across the
nine waste categories. The strongest classification performance was
observed for Vegetation, with an F1-score of $0.96$, followed by Cardboard
with an F1-score of $0.93$. The largest number of misclassifications
occurred between Plastic and Metal, where 8 Plastic samples were
classified as Metal and 7 Metal samples were classified as Plastic.
Miscellaneous Trash also showed some confusion with Plastic and Textile
Trash.

\subsubsection{MobileNetV2}

MobileNetV2 was evaluated using the same dataset split, preprocessing
procedure, and two-stage training strategy as EfficientNet-B0. The model
was loaded with its ImageNet pretrained weights, and its final 1000-class
classification layer was replaced with a new fully connected layer with
9 output neurons.

Initially, the feature extraction layers of MobileNetV2 were frozen,
while only the newly added classification layer was trained. The model
was trained for 10 epochs using the Adam optimizer with a learning rate
of $0.001$. This allowed the newly added classifier to adapt the
pretrained ImageNet features to the waste classification task without
modifying the pretrained feature extractor.

After the initial 10 epochs, the deeper layers of the network were
fine-tuned. Specifically, feature blocks 14--18 were unfrozen while the
earlier feature layers remained frozen. The learning rate was reduced to
$0.0001$ to allow the pretrained weights to be adjusted gradually without
significantly disrupting the previously learned features. The model was
subsequently fine-tuned for a further 5 epochs.

The best validation accuracy obtained during fine-tuning was $89.20\%$.
The final model was then evaluated on the previously unseen test set,
obtaining a test loss of $0.3444$ and a test accuracy of $88.50\%$.

\begin{table}[H]
    \centering
    \caption{MobileNetV2 classification report on the RealWaste test set.}
    \label{tab:mobilenetv2_classification}
    \begin{tabular}{lcccc}
        \hline
        Class & Precision & Recall & F1-score & Support \\
        \hline
        Cardboard & 0.94 & 0.93 & 0.93 & 69 \\
        Food Organics & 0.93 & 0.92 & 0.93 & 62 \\
        Glass & 0.85 & 0.89 & 0.87 & 63 \\
        Metal & 0.85 & 0.89 & 0.87 & 118 \\
        Miscellaneous Trash & 0.77 & 0.79 & 0.78 & 75 \\
        Paper & 0.99 & 0.88 & 0.93 & 75 \\
        Plastic & 0.87 & 0.84 & 0.86 & 138 \\
        Textile Trash & 0.88 & 0.96 & 0.92 & 48 \\
        Vegetation & 0.95 & 0.95 & 0.95 & 65 \\
        \hline
        Accuracy & & & 0.88 & 713 \\
        Macro Avg. & 0.89 & 0.89 & 0.89 & 713 \\
        Weighted Avg. & 0.89 & 0.88 & 0.89 & 713 \\
        \hline
    \end{tabular}
\end{table}

The confusion matrix for MobileNetV2 is shown in Table~\ref{tab:mobilenetv2_confusion}.

\begin{table}[H]
    \centering
    \caption{MobileNetV2 confusion matrix on the RealWaste test set.}
    \label{tab:mobilenetv2_confusion}
    \resizebox{\textwidth}{!}{
    \begin{tabular}{lccccccccc}
        \hline
        & Cardboard & Food Org. & Glass & Metal & Misc. Trash & Paper & Plastic & Textile & Vegetation \\
        \hline
        Cardboard
        & 64 & 0 & 2 & 2 & 0 & 1 & 0 & 0 & 0 \\
        Food Organics
        & 0 & 57 & 0 & 0 & 1 & 0 & 2 & 0 & 2 \\
        Glass
        & 1 & 0 & 56 & 3 & 2 & 0 & 1 & 0 & 0 \\
        Metal
        & 0 & 0 & 1 & 105 & 5 & 0 & 7 & 0 & 0 \\
        Misc. Trash
        & 0 & 2 & 1 & 1 & 59 & 0 & 6 & 5 & 1 \\
        Paper
        & 2 & 1 & 1 & 2 & 2 & 66 & 1 & 0 & 0 \\
        Plastic
        & 1 & 1 & 5 & 11 & 3 & 0 & 116 & 1 & 0 \\
        Textile Trash
        & 0 & 0 & 0 & 0 & 2 & 0 & 0 & 46 & 0 \\
        Vegetation
        & 0 & 0 & 0 & 0 & 3 & 0 & 0 & 0 & 62 \\
        \hline
    \end{tabular}
    }
\end{table}

The model correctly classified most samples in the Cardboard, Food
Organics, Paper, Textile Trash, and Vegetation categories. The most
noticeable confusion occurred between Plastic and Metal, where 11 Plastic
samples were classified as Metal and 7 Metal samples were classified as
Plastic. Miscellaneous Trash also showed relatively greater confusion with
Plastic and Textile Trash.

The test performance and model complexity of the two pre-trained
architectures are summarized in Tables~\ref{tab:pretrained_performance}
and~\ref{tab:model_complexity}.

\begin{table}[H]
    \centering
    \caption{Test performance comparison of the pre-trained lightweight architectures.}
    \label{tab:pretrained_performance}
    \begin{tabular}{lcc}
        \hline
        Model & Test Loss & Test Accuracy \\
        \hline
        EfficientNet-B0 & 0.3521 & 88.92\% \\
        MobileNetV2 & 0.3444 & 88.50\% \\
        \hline
    \end{tabular}
\end{table}

\begin{table}[H]
    \centering
    \caption{Model complexity of the pre-trained lightweight architectures.}
    \label{tab:model_complexity}
    \begin{tabular}{lcc}
        \hline
        Model & Total Parameters & Saved Model Size (MB) \\
        \hline
        EfficientNet-B0 & 4,019,077 & 15.62 \\
        MobileNetV2 & 2,235,401 & 8.76 \\
        \hline
    \end{tabular}
\end{table}

% ---------------------------------------------------------------------
\newpage
\subsection{Final Comparison \& Trade-off Discussion}

Table~\ref{tab:final_comparison} compares the custom resource-constrained
CNN (Model B) with the two fine-tuned lightweight SOTA models.
Computational cost is expressed as approximate multiply--accumulate
operations (MACs) per image. For Model B this was calculated from the
layer dimensions; for MobileNetV2 and EfficientNet-B0 the values are the
approximate figures reported for the original architectures at
$224\times224$ input, since replacing the final layer with a 9-class
layer changes the total only negligibly.

\begin{table}[H]
\centering
\caption{Final comparison of Model B and the fine-tuned SOTA models}
\label{tab:final_comparison}
\resizebox{\textwidth}{!}{
\begin{tabular}{lccc}
\hline
\textbf{Metric} & \textbf{Model B (custom)} & \textbf{MobileNetV2} & \textbf{EfficientNet-B0} \\
\hline
Initialization & Random (from scratch) & ImageNet pretrained & ImageNet pretrained \\
Input Resolution & $64\times64$ & $224\times224$ & $224\times224$ \\
Total Parameters & 47,210 & 2,235,401 & 4,019,077 \\
Parameters relative to Model B & $1\times$ & $47.4\times$ & $85.1\times$ \\
Model Size & $\approx$184\,KB (est.) & 8.76\,MB & 15.62\,MB \\
Approx. MACs per Image & $\approx$5.2\,M & $\approx$300\,M & $\approx$390\,M \\
Training Epochs & 20 & 15 & 20 \\
Test Accuracy & 63.11 $\pm$ 3.62\%$^{\dagger}$ & 88.50\% & 88.92\% \\
Weighted F1-score & 61.75 $\pm$ 3.77\%$^{\dagger}$ & 89\% & 89\% \\
\hline
\end{tabular}
}

\vspace{1mm}
\footnotesize $^{\dagger}$Mean $\pm$ SD over five independent runs. The
SOTA models were fine-tuned once.
\end{table}

\paragraph{Accuracy.}
The pre-trained models clearly outperformed the custom model. MobileNetV2
and EfficientNet-B0 reached 88.50\% and 88.92\% test accuracy, roughly 25
percentage points above Model B's mean of $63.11\pm3.62\%$. Even the best
of Model B's five runs (67.74\%) was more than 20 percentage points below
either SOTA model, so the gap is far larger than the variation between
runs. The SOTA results come from a single fine-tuning run each; however,
fine-tuning from fixed pretrained weights is generally much less
sensitive to the random seed than training from scratch. The gap was
largest on the difficult classes, Miscellaneous Trash and Textile Trash,
where the SOTA models reached recalls of 0.79--0.80 and 0.92--0.96
respectively, compared with Model B's five-run means of only 0.28 and
0.25. Three factors explain this. First, the
SOTA models start from features learned on over a million ImageNet
images, whereas Model B must learn everything from roughly 3,300
training images. Second, they have far greater capacity (deeper
inverted-residual blocks, squeeze-and-excitation in EfficientNet-B0).
Third, they operate on $224\times224$ inputs, which preserve fine texture
detail that is lost when images are reduced to $64\times64$ for Model B.
The accuracy comparison therefore reflects the whole design choice, not
architecture alone.

\paragraph{Memory footprint.}
Model B is the clear winner on memory. With 47,210 parameters its
weights require only about 184\,KB in 32-bit floating point (about 46\,KB
after 8-bit quantization), which fits within the flash memory of many
microcontrollers. MobileNetV2 requires 8.76\,MB and EfficientNet-B0
15.62\,MB, i.e.\ roughly 47--85 times more.  Activation memory during inference follows the same trend,
because the SOTA models process feature maps derived from $224\times224$
inputs rather than $64\times64$.

\paragraph{Computational cost.}
Model B requires roughly 5.2 million MACs per image, compared with
approximately 300 million for MobileNetV2 and 390 million for
EfficientNet-B0, a difference of about 58--75 times. This translates into
lower latency and energy per inference, which matters for battery-powered
or always-on devices such as a smart waste bin. Between the two SOTA
models, MobileNetV2 achieved almost the same accuracy as EfficientNet-B0
(only 0.42 percentage points lower) with about 55\% of the parameters and
fewer operations, making it the more efficient of the two for this task.

\paragraph{Advantages and limitations of a custom sub-100k model.}
Building a custom model from scratch gives full control over every layer,
so it can be sized exactly to a hardware budget, uses only simple
operations (depthwise and pointwise convolutions, ReLU, pooling) that are
widely supported by embedded inference frameworks, and has no dependence
on large pretrained checkpoints. Its main limitations are limited
representational capacity and the lack of pretrained knowledge, which
together make it data-hungry and weaker on visually diverse classes; it
also requires more design and tuning effort to reach acceptable accuracy.

\paragraph{Advantages and limitations of pre-trained lightweight models.}
Industry-standard pre-trained models deliver high accuracy quickly, with
only 15--20 epochs of fine-tuning, and benefit from well-tested
architectures and tooling. Their limitations are a footprint measured in
megabytes and a compute cost in the hundreds of millions of operations
per image, which rules out the smallest edge devices and increases
latency and energy use on modest hardware.

\paragraph{Conclusion.}
If the deployment target can accommodate a few megabytes of memory, such
as a smartphone or an embedded Linux board, a fine-tuned MobileNetV2 offers
the best balance of accuracy and efficiency in this study. If the target
is a microcontroller with only hundreds of kilobytes of memory, Model B
is one of the few viable options, at the cost of substantially lower
accuracy. Ways to reduce this gap are discussed in
Section~\ref{sec:future_work}.

% ---------------------------------------------------------------------
\subsection{Future Work}
\label{sec:future_work}

\begin{itemize}
    \item Train all models, including the SOTA models, over the same set of fixed random seeds for full reproducibility.
    \item Seed all sources of randomness and enable deterministic training.
    \item Train Model B for more epochs with a learning-rate schedule and early stopping.
  
\end{itemize}

% End of document
\end{document}